\section{Justificativa}
O envelhecimento cutâneo envolve alterações na matriz extracelular, na hidratação e na pigmentação. Investigar moléculas que interajam com proteínas ligadas a esses processos pode orientar estudos sobre atividade enzimática e aplicações cutâneas \parencite{liu2023,papaccio2022}. A bioprospecção de microalgas amplia a busca por substâncias bioativas de interesse biotecnológico \parencite{khemiri2023}.

A pesquisa também se insere na bioeconomia azul, que promove conhecimento e desenvolvimento tecnológico a partir de recursos marinhos. \textcite{pereira2019} documentaram a produção industrial de biomassa de \textit{Tetraselmis} sp. CTP4. Para \textit{T. chuii}, estudos avaliaram compostos neuroprotetores em extratos e a atividade antioxidante da biomassa após digestão simulada; a biomassa seca da espécie também recebeu autorização para usos alimentares especificados \parencite{cokdinleyen2024,moon2023,aecosan2014tetraselmis}. Essas pesquisas fundamentam o interesse em ampliar a investigação de sua composição.

A SPE é relevante por sua fração polissacarídica e pelas possíveis interações com proteínas. Polissacarídeos sulfatados de \textit{Chlamydomonas reinhardtii} apresentaram atividade antioxidante e efeitos sobre enzimas associadas ao envelhecimento cutâneo \parencite{falcao2020}. Estudos com vesículas extracelulares de \textit{T. chuii} relataram biocompatibilidade e propriedades antioxidantes e anti-inflamatórias em modelos celulares da pele \parencite{adamo2024,gargano2026}. Essas evidências motivam a investigação da SPE de \textit{T. chuii} e de suas possíveis aplicações cutâneas.

\Needspace{8\baselineskip}
A seleção dos seis alvos humanos contempla mecanismos distintos associados à pele. MMP-1, MMP-2, MMP-9 e elastase participam da remodelação da matriz extracelular; a hialuronidase 1 atua sobre o ácido hialurônico, relacionado à hidratação; e TYRP1 está associada à melanogênese e à pigmentação \parencite{chao2007,lai2017,liu2025,macdonald2002}. A comparação desses alvos organiza a análise por processo biológico e receptor. Em cada proteína, os ligantes comerciais permitem contextualizar os resultados; quando disponíveis, os ligantes coocristalizados oferecem uma referência estrutural para avaliar a posição e os contatos previstos.

O docking molecular permite realizar uma triagem virtual dos ligantes nos sítios de ligação dos receptores. O algoritmo explora posições e orientações possíveis, atribui pontuações às poses e, com a análise do PLIP, identifica contatos entre ligantes e resíduos proteicos. A avaliação conjunta desses dados permite classificar os compostos avaliados em cada receptor e priorizar moléculas para estudos experimentais de atividade enzimática, segurança e permeação cutânea \parencite{salentin2015plip,trott2010,warren2006}. Assim, os resultados computacionais orientam o planejamento das etapas experimentais subsequentes.

O estudo também se relaciona aos Objetivos de Desenvolvimento Sustentável (ODS) 9 e 14, referentes à inovação e à infraestrutura e à conservação e ao uso sustentável dos oceanos \parencite{onu2015agenda}. Essa relação se estabelece pela produção de conhecimento científico sobre recursos microalgais e converge com as metas 9.5, de fortalecimento da pesquisa e da capacidade tecnológica, e 14.a, de ampliação do conhecimento científico e da tecnologia marinha. A prospecção molecular desenvolvida neste trabalho contribui para esse campo de investigação.
